\bookchapter{第十章：流動支付與生活服務}

\chapterintro{講真，我爸媽嗰一輩人對「手機畀錢」呢件事，一開始係抗拒嘅。「錢喺手機入面，被人偷咗點算？」呢個擔心好正常，但現實係：你褸袋入面揣住嘅現金，遺失嘅概率比手機支付被盜刷嘅概率高得多。手機支付唔係將錢存入手機，佢淨係一把鑰匙，一把打開你銀行賬戶嘅鑰匙。}

\booksection{一、到底點樣綁定銀行卡？}

唔管你係用支付寶定微信，第一次綁卡都要行呢幾步：

\begin{enumerate}
  \item 打開App，揾到「我嘅」或「錢包」。點「添加銀行卡」。
  \item 輸入卡號。而家大部分App都支援直接影相識別卡號，唔使一個個數字敲。
  \item 銀行會發一條驗證碼短訊到你預留嘅手機號上。填入去，設定一個6位數字嘅支付密碼。
  \item 呢個支付密碼同你嘅銀行卡密碼、鎖屏密碼，最好係三個唔同嘅密碼，唔好圖省事設成一樣。
\end{enumerate}

有個小竅門：如果擔心安全問題，可以去銀行App入面開一張II類或III類電子賬戶，專門用嚟綁定手機支付。呢張卡入面淨係放平時零花嘅錢，就算出咗事，損失都有限。

\booksection{二、付款碼係「邊個睇邊個收錢」}

好多人搞唔清楚付款碼同收款碼嘅區別，呢個區別好重要。

\begin{itemize}
  \item \textbf{付款碼}：你展示畀收銀員掃嗰個條形碼+二維碼。呢個碼相當於你嘅「電子銀包」，邊個掃咗邊個就可以由你賬上掃走錢。所以\textbf{千萬唔好截圖發畀任何人}。呢個碼每分鐘自動刷新一次，就係為咗防止被盜用。
  \item \textbf{收款碼}：別人掃咗呢個碼畀你轉賬。呢個碼可以截圖保存，可以發畀別人，安全。
\end{itemize}

認準呢個區別，可以避開好多詐騙陷阱。凡係有人等你「將付款碼截圖發過嚟」嘅，百分之百係騙子。

\booksection{三、呢啲安全設定，花5分鐘設好可以慳好多心}

\begin{itemize}
  \item \textbf{指紋/面容支付}：喺支付寶同微信嘅支付設定入面，將生物支付打開。咁樣每次付款要按指紋或掃臉，就算手機丟咗，執到嘅人都冇法使你嘅錢。
  \item \textbf{安全鎖}：支付寶同微信都有「安全鎖」功能。打開之後，進入App嘅「錢包」頁面需要額外驗證一次身份。路徑：支付寶 我嘅 設定 帳號與安全 安全中心 安全鎖。
  \item \textbf{免密支付檢查}：好多人簽約咗各種免密支付，自己都唔記得。去支付寶同微信嘅「支付設定」入面，揾到「免密支付/自動扣款」，睇下簽約咗邊啲服務，將唔需要嘅關掉。
  \item \textbf{登錄設備管理}：定期睇下邊啲設備登錄過你嘅帳號。如果發現唔識嘅設備，即刻踢走。路徑：支付寶 設定 帳號與安全 登錄設備管理。
\end{itemize}

\booksection{四、公交地鐵：手機就係交通卡}

而家出門坐公交地鐵，已經好少人用實體嘅交通卡喇。手機刷閘機，有兩條路可以行：

\begin{itemize}
  \item \textbf{手機錢包入面嘅虛擬交通卡}：喺蘋果錢包、華為錢包、小米錢包入面直接開通一張虛擬交通卡。呢個方案嘅好處係：刷嗰陣唔使解鎖手機，直接靠上去就得，響應速度最快，而且冇網都可以刷。部分城市嘅虛擬交通卡仲有同實體卡一樣嘅換乘優惠。
  \item \textbf{支付寶/微信嘅乘車碼}：喺App入面搜尋「乘車碼」，開通你所在城市嘅公交或地鐵乘車碼。呢個方案嘅好處係：唔使充值，直接由餘額或綁定嘅銀行卡入面扣。但需要聯網，而且刷之前要打開App揾到碼。
\end{itemize}

我個人更推薦用手機錢包嘅虛擬交通卡。雖然第一次要充個10蚊20蚊入去，但用起嚟真係方便太多，尤其係趕地鐵嗰陣，唔使停下嚟解鎖手機揾App。

\booksection{五、生活繳費、外賣、打車}

\begin{itemize}
  \item \textbf{水電氣繳費}：喺支付寶或微信嘅「生活繳費」入面，將你屋企嘅水、電、燃氣戶號綁定好。以後每個月收到賬單，直接點入去就可以交。戶號一般喺你屋企嘅紙質賬單上可以搵到，或者打客服電話問。
  \item \textbf{外賣比價}：同一間舖喺美團同餓了麼上嘅價格同優惠經常唔一樣。落單前兩邊都睇一眼，有時候可以慳幾蚊。另外，美團嗰個「拼好飯」功能，同朋友一齊點，價格可以平唔少。
  \item \textbf{打車安全}：喺滴滴或者高德打車入面，記得將「緊急聯絡人」同「行程分享」設定好。打車嗰陣，App會自動將你嘅行程分享畀你設定嘅聯絡人。另外，深夜打車優先選「的士」而唔係「快車」，的士有行業監管，安全系數更高啲。
\end{itemize}

\booksection{六、用手機搞掂醫保、社保、政務}

呢幾年嘅變化好大，以前要跑大廳排隊嘅好多事，而家手機上幾分鐘就辦完。

\begin{itemize}
  \item \textbf{醫保電子憑證}：喺支付寶或微信入面搜「醫保電子憑證」，激活之後，去醫院睇病、藥房買藥直接刷手機就得，唔使帶實體醫保卡。第一次需要人臉識別驗證，之後每次用嗰陣會生成一個動態碼。
  \item \textbf{社保查詢}：下載「掌上12333」App或者用支付寶入面嘅「市民中心」，可以查社保繳納記錄、養老金賬戶餘額、失業金申領等。
  \item \textbf{個税}：每年3月到6月嘅個税匯算清繳，喺「個人所得税」App上就可以搞掂。記得睇下有冇可以申報嘅專項附加扣除（子女教育、房貸利息、贍養老人等），呢啲都可以幫你少交税或者退税。
\end{itemize}

\begin{tcbwarn}{記住呢幾條，防騙保平安}
\begin{itemize}
  \item 公檢法、銀行、客服，唔會等你轉賬到「安全賬戶」。
  \item 任何人等你將付款碼截圖發過去，都係騙子。
  \item 陌生連結唔好點，尤其係「ETC過期」「社保卡停用」呢類短訊。
  \item 網購退款會原路返回，唔需要你提供銀行卡號。
\end{itemize}

\end{tcbwarn}

\begin{tcbpractice}{而家就做：畀你嘅支付App做一次安全體檢}
\begin{enumerate}
  \item 打開支付寶同微信，確認已開啟指紋/面容支付。
  \item 揾到「安全鎖」功能並打開。
  \item 檢查「免密支付/自動扣款」列表，關掉唔再需要嘅服務。
  \item 查看「登錄設備」列表，移除唔識嘅設備。
  \item 激活你嘅醫保電子憑證，下次去醫院就唔使帶實體卡喇。
\end{enumerate}

\end{tcbpractice}
